% ============================================================
%  06_chapitre1_presentation.tex  –  Chapitre 1 : Contexte Général
% ============================================================
\chapter{Contexte Général du Projet}

\section*{Introduction}
\addcontentsline{toc}{section}{Introduction}

Ce premier chapitre pose les bases de notre travail. Il présente l'organisation hôte,
le problème auquel notre solution répond, l'étude de la situation existante, la solution
adoptée et la méthodologie de travail choisie.

% ------------------------------------------------------------
\section{Cadre du Projet}

\subsection{Présentation de l'organisme d'accueil : Whitecape Technologies}

Whitecape Technologies est une entreprise franco-tunisienne spécialisée dans le
développement de solutions informatiques sur mesure. Située à Sousse, en Tunisie,
l'entreprise possède une vaste expérience en gestion de projets tant localement qu'en
modèle nearshore. Whitecape Technologies se distingue par son expertise dans la
création de solutions innovantes adaptées aux besoins spécifiques de chaque client,
en s'appuyant sur des équipes pluridisciplinaires et hautement qualifiées.

\begin{figure}[H]
  \centering
  
\includegraphics[width=0.35\textwidth]{whitecape_technologies_logo.jpg}
  \caption{Logo de Whitecape Technologies}
  \label{fig:logo_whitecape}
\end{figure}

\subsection{Objectif et Contexte}

Fondée en 2007, Whitecape Technologies vise à devenir un leader dans le secteur
des technologies de l'information en Tunisie et au-delà. L'entreprise se concentre
sur la fourniture de solutions technologiques avancées pour aider les entreprises à
optimiser leurs opérations, améliorer leur productivité et atteindre leurs objectifs
stratégiques.

\subsection{Services et Solutions}

\paragraph{Développement et ingénierie}
\begin{itemize}
  \item Développement d'applications web dynamiques et interactives.
  \item Développement d'applications mobiles iOS et Android.
  \item Solutions sur mesure et DevOps / automatisation (\ac{CI}/\ac{CD}).
\end{itemize}

\paragraph{Consultation en technologie}
Audit, évaluation et définition de la feuille de route technologique.

\paragraph{Intégration des systèmes}
Intégration ERP et CRM, support technique, maintenance préventive et corrective.

\paragraph{Services d'Assurance Qualité}
Tests manuels et automatisés, frameworks d'automatisation, audits qualité.

\medskip
\noindent\textbf{Fiche technique :}
\begin{itemize}
  \item \textbf{Entreprise :} Whitecape Technologies
  \item \textbf{Email :} \href{mailto:info@whitecapetech.com}{info@whitecapetech.com}
  \item \textbf{Téléphone :} +216 31 20 88 88 – +33 5 32 10 80 56
  \item \textbf{Adresse :} Avenue de la République, Akouda 4022, Sousse – Tunisie
  \item \textbf{Site web :} \url{https://www.whitecapetech.com}
\end{itemize}

\subsection{Problématique}

Le marché de la formation en ligne connaît une croissance exponentielle. Cependant,
les solutions existantes présentent des lacunes majeures~:

\begin{itemize}
  \item \textbf{Absence de personnalisation} des parcours selon le niveau et le rythme
        de chaque apprenant.
  \item \textbf{Évaluations statiques} sans feedback pédagogique contextuel.
  \item \textbf{Absence d'assistance 24h/24} pour accompagner l'étudiant.
  \item \textbf{Taux d'abandon élevé} dû au manque de motivation.
  \item \textbf{Gestion administrative lourde} pour enseignants et administrateurs.
\end{itemize}

\medskip
La question centrale est donc~:

\begin{quote}
\textit{Comment concevoir une plateforme e-learning moderne, scalable et intelligente,
qui exploite l'\ac{IA} générative pour personnaliser l'expérience de chaque apprenant
tout en simplifiant la gestion pour les enseignants et administrateurs~?}
\end{quote}

\subsection{Étude de l'existant}

\begin{table}[H]
  \centering
  \caption{Analyse comparative des plateformes e-learning existantes}
  \label{tab:existant}
  \begin{tabularx}{\textwidth}{lXXc}
    \toprule
    \textbf{Plateforme} & \textbf{Points forts} & \textbf{Limitations} & \textbf{\ac{IA}~?} \\
    \midrule
    Moodle      & Open source, flexible, large communauté & Interface vieillissante, installation complexe & Non \\
    \midrule
    Udemy       & Large catalogue (200k+ cours), mobile  & Pas de personnalisation du parcours & Partiel \\
    \midrule
    Coursera    & Certifications reconnues, partenariats universités & Coûteux, peu adapté PME tunisiennes & Partiel \\
    \midrule
    Google Classroom & Intégration Google, simplicité & Très limité, pas de marketplace & Non \\
    \midrule
    \textbf{EduBoard} & IA complète, multi-rôles, chapitres verrouillés, certificats auto & Solution en développement & \textbf{Oui (complet)} \\
    \bottomrule
  \end{tabularx}
\end{table}

\subsection{Solution proposée : EduBoard}

EduBoard est une plateforme e-learning intelligente et complète, développée avec des
technologies modernes. Elle répond à l'ensemble des lacunes identifiées~:

\begin{enumerate}
  \item Gestion multi-rôles avec tableau de bord dédié (Admin, Enseignant, Étudiant).
  \item Cours structurés en chapitres et leçons (article, PDF).
  \item Chapitres verrouillés séquentiellement : accès au chapitre $N$ conditionné
        par la complétion du chapitre $N-1$.
  \item Quiz adaptatifs avec feedback \ac{IA} généré par Claude (Anthropic).
  \item Chatbot pédagogique 7j/7 propulsé par l'API Groq.
  \item Moteur de recommandation basé sur le comportement et la progression.
  \item Génération automatique de certificats PDF via QuestPDF avec QR code.
  \item Système de paiement Flouci pour les cours payants.
\end{enumerate}

% ------------------------------------------------------------
\section{Méthodologie de Travail : SCRUM}

\subsection{Présentation du Framework SCRUM}

\noindent Conformément aux pratiques de gestion de projet adoptées par \textbf{Whitecape Technologies}, le projet EduBoard a été conduit selon une démarche agile inspirée du framework \textbf{SCRUM}. Ce choix se justifie par la nature intrinsèquement évolutive des besoins d'une plateforme e-learning, la nécessité de produire des livrables testables à chaque itération et l'importance d'une communication régulière entre les parties prenantes.

\vspace{0.5cm}

\noindent SCRUM est un cadre de travail agile léger mais puissant, issu de la métaphore du rugby (la mêlée, ou \textit{scrum} en anglais). Il repose sur trois piliers fondamentaux : la \textbf{transparence} (tous les aspects du processus sont visibles par tous), l'\textbf{inspection} (les artefacts et la progression sont régulièrement examinés) et l'\textbf{adaptation} (le plan est ajusté en fonction des résultats des inspections).

\vspace{0.5cm}

\noindent Les éléments constitutifs de SCRUM se déclinent en trois catégories :

\textbf{Rôles :}
\begin{itemize}
  \item \textbf{Product Owner :} Représente les intérêts des parties prenantes, gère et priorise le Product Backlog. Dans notre contexte : l'encadrant professionnel Wael Boukadida.
  \item \textbf{Scrum Master :} Garant de l'application du framework, facilitateur et protecteur de l'équipe. Rôle assumé par l'auteure.
  \item \textbf{Équipe de Développement :} Auto-organisée et pluridisciplinaire. Composée de l'auteure pour le développement full-stack.
\end{itemize}

\textbf{Artefacts :}
\begin{itemize}
  \item \textbf{Product Backlog :} Liste priorisée de toutes les fonctionnalités souhaitées pour le produit final.
  \item \textbf{Sprint Backlog :} Sous-ensemble du Product Backlog sélectionné pour le sprint en cours.
  \item \textbf{Incrément :} Somme de toutes les user stories complétées, constituant une version potentiellement livrable du produit.
\end{itemize}

\textbf{Cérémonies :}
\begin{itemize}
  \item \textbf{Sprint Planning :} Réunion de planification du sprint (objectifs, user stories sélectionnées, estimations).
  \item \textbf{Daily Scrum :} Point quotidien de 15 minutes (avancement, blocages, plan du jour).
  \item \textbf{Sprint Review :} Démonstration et validation des fonctionnalités développées.
  \item \textbf{Sprint Retrospective :} Analyse du processus pour améliorer la productivité et la qualité.
  \vspace{0.7cm}
\end{itemize}

\begin{figure}[H]
    
  \centering
  \includegraphics[width=0.85\textwidth]
  {scrum.png}
  \caption{Cycle de vie du framework SCRUM}
  \label{fig:scrum}
\end{figure}


\section*{Conclusion}
\addcontentsline{toc}{section}{Conclusion}

Ce chapitre a posé le cadre du projet~: présentation de l'organisme d'accueil,
analyse de la problématique et des solutions existantes, définition de la solution
EduBoard et planification SCRUM. Le chapitre suivant formalise les besoins fonctionnels
et non fonctionnels de la plateforme.
